\chapter*{ABSTRACT}
\addcontentsline{toc}{chapter}{ABSTRACT}

This thesis, \emph{\ThesisTitle}, studies longitudinal multimodal survival prediction of MCI progression, using the CNN--T-LSTM framework of Aghajanian et al. as its baseline. MRI, concatenation and pairwise fusion models are evaluated before developing a compact CP architecture.

The frozen ADNI cohort contains 359 participants and 106 events, with three MRI visits within an 18-month eligibility window. The third MRI defines the prediction reference time; events are ascertained after the eligibility landmark using MMSE/CDR thresholds. Retrospective clinical alignment within $\pm90$ days of each MRI limits claims about prospective input availability.

MRI--Radiomics (MR), MRI--Clinical (MC) and Radiomics--Clinical (RC) interactions initially use outer products and dynamic gated residual routing. CP factorizes the bilinear \emph{weight tensor}, retaining the subsequent MLP. The three-pair interaction core falls from 229,696 to 3,712 parameters, while three-seed mean validation C-index changes from 0.8344 to 0.8278. Ablation motivates RC-Free Pairwise Fusion, which retains all modalities and only MR/MC. Its four routes use input-dependent sigmoid gates and learnable global scales, followed by T-LSTM, temporal attention and a Cox head.

RC-Free achieves standardized validation C-index $0.8347\pm0.0237$ (sample SD over three seeds). At seed 20260727, it obtains the highest ablation validation score of 0.8417 and a post-hoc test score of 0.8750; the local CNN--PCA--T-LSTM Adaptation obtains 0.6625 and 0.5379, respectively. RC-Free is selected by validation and structural simplicity, although full CP has a higher post-hoc test score. Statistical superiority is not established. Historical test access makes these test results post-hoc; independent final evaluation requires an independent cohort or a previously unused held-out dataset.

\textbf{Keywords:} Alzheimer's disease, MCI, longitudinal MRI, radiomics, tensor fusion, CP factorization, gated residual routing, T-LSTM, survival analysis.
